%!TEX root = ../thesis.tex

\chapter{Introduction}

\section{Physical Applications and Reproducibility}
% <<SECTION_NEW_BEGIN>>
Compressible-Euler and HRSC methods underpin engineering and scientific calculations across several settings. External aerodynamics use them for supersonic intakes and hypersonic re-entry, while turbomachinery analysis uses them to capture shocks inside transonic cascades. Inertial-confinement fusion needs them at imploding interfaces~\citep{atzeni_meyer_2004}; plasma and magnetohydrodynamic calculations sit on the same finite-volume core~\citep{bard_dorelli_2014}. These users run design loops or operating-envelope sweeps that need fast, accurate solutions.

%

Implementation choices vary between groups, hardware vendors, and compiler toolchains. Two builds of the same algorithm can disagree from precision, FMA contraction, reduction ordering, or branch-rule choices near zero wave speeds. Report 1 quantifies that variability on a controlled Euler validation suite.
% <<SECTION_NEW_END>>

\section{HRSC Schemes for Discontinuous Compressible Flows}
% <<SECTION_1_BEGIN>>
Compressible HRSC methods evolve conserved cell averages with Riemann-solver fluxes, capturing shocks and contact discontinuities without tracking them explicitly~\citep{toro2009}. The same features that make these methods useful also make their precision and backend reproducibility worth measuring.
% <<SECTION_1_END>>

\section{Precision and Hardware Reproducibility Problem}
% <<SECTION_2_BEGIN>>
For such schemes, the finite-volume algorithm alone does not guarantee reproducibility. Two runs can share reconstruction, Riemann solver, and time step, yet differ because binary32 and binary64 round differently, compilers fuse or reassociate expressions, or device arithmetic and parallel ordering change rounded operations \citep{goldberg_1991,ieee754_2019}. These mechanisms matter because shock-capturing updates contain branch thresholds, limiter switching, and cancellation near discontinuities. Recent CFD studies treat reduced or mixed precision as case-dependent: Brogi et al. report OpenFOAM evidence across CFD workloads, while Wang, Xia, and Chen study heterogeneous hybrid precision for compressible flow. That work sits against the wider mixed-precision literature surveyed by \citet{higham_mary_2022} \citep{brogi_etal_2024,wang_xia_chen_2025}. These papers motivate validation here; they do not establish fp32 suitability. CUDA can differ from CPU backends and across GPU devices, so this report measures rather than assumes hardware reproducibility.
% <<SECTION_2_END>>

\section{Report Scope and Contribution}
% <<SECTION_3_BEGIN>>
This report measures how precision and hardware/backend choices affect a
Riemann-solver-based HRSC method on a controlled Euler validation suite. MHD
motivates the longer project, but its magnetic-field constraint and
divergence-control choices require separate validation before they can support
conclusions. Within the tested Euler cases, the report contributes validation
evidence for the stand-alone HRSC solver, quantified fp32/fp64 and matched
CPU/GPU saved-output comparisons, and measured sensitivity to selected compiler,
solver, and branch-rule variations.
% <<SECTION_3_END>>

\section{Report Structure}
% <<SECTION_4_BEGIN>>
Chapters~2--4 establish background, method, and design; Chapters~5--6 present results and discussion; Chapter~7 concludes.
% <<SECTION_4_END>>
